\section{主结果与两层几何}
\label{sec:main-results}

__CONSTANTS__
此处的下确界针对定义\ref{def:admissible}中的全部可测可行控制，不预先限制切换次数。

\subsection{约束可行性与成本最优性的区分}
__PEAK__
先从能否保持感染阈值出发定义第一层集合。
记 $\mathcal C_{\bar q}$ 为所有可测的 $[0,\bar q]$ 值控制，$0\le\bar q\le1$；
定义该控制类时不预先要求轨道满足 $i\le K$。令
\begin{align}
 \Acal&=\{x\in\Dcal:P_0(x)\le K\},\label{eq:overview-safe}\\
 \mathcal A^0&=\{x\in\Dcal:i_x^q(t)\le K\ \text{对所有 }t\ge0\text{ 及 }q\in\mathcal C_1\},
 \label{eq:all-control-safe}\\
 \mathcal V_{\bar q}&=\{x\in\Dcal:\text{存在 }q\in\mathcal C_{\bar q},\ i_x^q(t)\le K\ \text{对所有 }t\ge0\}.
 \label{eq:viability-definition}
\end{align}
三者分别表示零控制安全集、所有允许控制均安全的初值集和存在可行控制的初值集。
这种先判定可行性、再比较成本的层级已有研究先例\citep{AvramFreddiGoreac2022,FreddiGoreac2024}。
本模型满足
\begin{equation}
 \mathcal A^0=\Acal=\mathcal V_0,\qquad \mathcal V_1=\Dcal,
 \label{eq:viability-endpoints}
\end{equation}
证明见引理\ref{lem:safe-invariant}；受限能力 $\bar q<1$ 的可行包络见命题\ref{prop:restricted-feasibility}。
因此，完整控制区间下所有物理初值均可行，但可行性本身尚未决定何时付出控制成本。
第二层才是按最小成本确定的等待区、完全跟踪区和容量平台。

\subsection{全域反馈与阶段结构}
\begin{theorem}[常接触率下的全域最优反馈与唯一性]
\label{thm:main-overview}
在式\eqref{eq:constants}的条件下，对每个 $x_0\in\Dcal$，存在有限成本最优控制。
定义\ref{def:candidate}中的分区反馈 $Q^*$ 实现该最优值，其候选成本满足 $U(x_0)=V(x_0)$，
比较范围为全部可测可行控制。
从同一初值出发的最优状态轨道唯一，最优控制在 $[0,\infty)$ 上几乎处处唯一。
不安全初值的阶段顺序为 $0\to1\to0$ 或 $0\to q_B(s)\to1\to0$，允许阶段长度为零，
其中 $q_B(s)=1-h/s$。两类策略均在首次进入 $\Acal$ 后永久取 $q=0$。
辅助条带中的非平凡切换曲线 $\Gamma_K$ 及其唯一容量交点 $(s_B,K)$ 由命题\ref{prop:Gamma-geometry}、推论\ref{cor:Gamma-labels}刻画。
若最优轨道含有正长度容量阶段，则在 $s_B>h$ 处转入完全跟踪，而不继续维持感染阈值至 $s=h$。
\end{theorem}
\begin{proof}
这是后文定理\ref{thm:verification}、定理\ref{thm:unique-q}和推论\ref{cor:optimal-synthesis}的合并陈述。
后文依次构造几何与候选成本、验证其最优性、再证明唯一性，均不预先使用本概览结论。
\end{proof}

表\ref{tab:optimal-stages}列出阶段顺序；分区的精确定义见第\ref{sec:candidate}节。
表中集合均与物理域 $\Dcal$ 取交，超出物理三角形的辅助切换点不对应可实现的阶段。
\begin{table}[H]
\centering
\caption{完整控制区间下的最优阶段。$0$、$1$ 和 $q_B$ 分别表示零控制、完全跟踪与容量切向控制。}
\label{tab:optimal-stages}
\begin{tabular}{@{}lll@{}}
\toprule
初值区域 & 阶段顺序 & 停止付费的状态\\
\midrule
$\Acal$ & 永久 $0$ & 初始已安全\\
$\mathcal T\cup\Gamma_K$ & $1\to0$ & 首次进入 $\Acal$\\
$\mathcal W_\Gamma$ & $0\to1\to0$ & 首次进入 $\Acal$\\
$\mathcal W_K\cup\mathcal B$ & $0\to q_B\to1\to0$ & 首次进入 $\Acal$\\
\bottomrule
\end{tabular}
\end{table}
在 $\mathcal B$ 出发时等待段为零；临界自然轨道 $\Phi=\Phi_B$ 对应零长度容量段。
图\ref{fig:two-layer-regions}将约束可行性与成本最优性分别画出。
\begin{figure}[H]
\centering
\includegraphics[width=\linewidth]{numerical_scenarios/FigN1_two_layers.pdf}
\caption{两层几何。左图为 $\mathcal A^0=\Acal$、受限能力的 $\mathcal V_{\bar q}$ 与 $\mathcal V_1=\Dcal$；
右图为 $q\in[0,1]$ 下的成本最优分区，图中曲线均按解析公式计算。
参数为 $\beta=0.5,c=2,\gamma=0.3,K=0.15$，左图另取 $\bar q=0.4$。
点线 $s=h$ 为自然感染增长阈值，虚线 $s+i=1$ 为物理总量边界。}
\label{fig:two-layer-regions}
\end{figure}
